%%-----------------------------------------------------
%%   And the last abstract from Rabat. The e-mail of
%%authors is:  mikram@fsr.ac.ma
%%----------------------------------------------------
%% This document created by Scientific Word (R)
%% Version 2.0


\documentclass[11pt,thmsa]{article}
\newcommand{\R}{\mbox{\rm\bf R}} % the field of reals
%\usepackage{amsfonts}
%%%%%
%\usepackage{sw20lart}

%TCIDATA{TCIstyle=Article/art4.lat,lart,article}

%\input tcilatex
\begin{document}


\begin{center}
\textbf{Computation of Normal Forms of Hamiltonian Systems}

\textbf{in the Presence of Poisson Commuting integrals}

(Liouville-Integrability and Birkhoff Normal Forms)\\[0.2cm]
\end{center}

\begin{center}
\textit{J. Mikram \& F. Zinoun}

\textit{D\'epartement de Math\'ematiques et Informatique}

\textit{Facult\'e des Sciences de Rabat, Morocco}\\[1cm]
\end{center}
Consider a Hamiltonian system with $n$ degrees of freedom having the
origin as an equilibrium point, and assume that the Hamiltonian function $%
H(z)$ is of class $C^\infty $ in a neighborhood of the origin. By setting $%
H(0)=0$, the Hamiltonian function can be written as 
\[
H=\frac 12\left\langle Sz,z\right\rangle +O\left( \left| z\right| ^3\right) 
\]
and the corresponding Hamiltonian system is 
\begin{equation}
\frac{dz}{dt}=JH_z=JS_z+O\left( \left| z\right| ^2\right) ,\;\;\;J=\left( 
\begin{array}{cc}
0 & I \\ 
-I & 0
\end{array}
\right) 
\end{equation}
where $S$ is the Hessian matrix of $H$ at the origin, $\left\langle
.,.\right\rangle $ denotes the Euclidean inner product of {\R}$^{2n}$, $I$
is the identity matrix of degree $n$ and $O\left( \left| z\right| ^d\right) $
is a function whose Taylor expansion at the origin begins with terms of
order $\geq d$. If the origin is an elliptic equilibrium point, we can find
a (real) linear canonical transformation so that the Hamiltonian function is
of the form 
\[
H=\frac 12\sum\limits_{j=1}^n\lambda _j\left( x_j^2+y_j^2\right) +O\left(
\left| z\right| ^3\right) ,\;\;\lambda _j\in \R\
\]
A preliminary reduction to normal form consists in carrying out a canonical
change in variables allowing to pass from the Hamiltonian $H$ to a simpler
(normal form) $\widetilde{H}$, satisfying 
\[
adH_2\left( \widetilde{H}\right) =0
\]
$adH_2$ being the linear operator $\left\{ H_2,-\right\} $ where $H_2=\frac
12\sum\limits_{j=1}^n\lambda _j\left( x_j^2+y_j^2\right) $ and the brackets
denote the canonical Poisson bracket 
\[
\left\{ F,G\right\} =\sum\limits_{j=1}^n\left( \frac{\partial F}{\partial
x_j}\frac{\partial G}{\partial y_j}-\frac{\partial F}{\partial y_j}\frac{%
\partial G}{\partial x_j}\right) 
\]
The aim of this talk is to show how one can take profitably the
existence of first integrals (i.e., quantities $F$ that are preserved by the
flow of the Hamiltonian system, or equivalently $\left\{ H,F\right\} =0$) to
produce a simplified normal form.\\ The motivation for considering the first
integrals of the Hamiltonian (1) in view of a further reduction of the
normal form comes from the following fact:\\ Let $F$ be a first integral of
(1), that is 
\[
\left\{ H,F\right\} =0
\]
Suppose we have taken $H$ into a normal form $\widetilde{H}$ via a canonical
coordinate transformation $T$. Then, if we consider the transformation of $F$
under $T$, writing $F\circ T\equiv \widetilde{F}$ (not necessarily in normal
form!), we have 
\[
\left\{ \widetilde{H},\widetilde{F}\right\} =0
\]
because $T$ is canonical and then it preserves the Poisson bracket.
Therefore, the presence of a first integral for (1) will pose some
restriction to the normal form $\widetilde{H}$. In fact, while in general $%
\widetilde{H}$ only satisfies (with obvious notation) 
\[
\widetilde{H}\in Ker\left( adH_2\right) \;\;,
\]
in the presence of the first integral we will also have 
\[
\widetilde{H}\in Ker\left( ad\widetilde{F}\right) 
\]
which can lead to a simplification of the normal form.\\[0.2cm] This
obvious, but important, fact prompts to examine in a more general context
for general Hamiltonian structures the role of first integrals in further
reducing the normal form unfolding of Hamiltonian systems. The result is
that, via a normal form approach, we can reduce simultaneously and in an
algorithmic way the Hamiltonian together with its Poisson commuting
integrals into a joint normal form which is given a simple characterization.\\[0.2cm] A prominent class of Hamiltonian systems for illustrating the
joint normal form approach are Liouville-integrable systems, i.e., roughly, $%
n$ degrees of freedom Hamiltonians with $n$ functionally independent and
Poisson commuting integrals. We have shown that a number of interesting
results on computation of Birkhoff normal forms for Liouville-integrable
Hamiltonian systems can be seen as a direct consequence of the joint normal
form approach.\\[0.2cm] As a Computer Algebra Application, we choose an
integrable case of the generalized H\'enon-Heiles system. This provides an
interesting example to show the importance of the joint normal form
algorithm in taking to Birkhoff normal form (and then solving)
Liouville-integrable Hamiltonian systems.\\[0.2cm] This work is an extension
(and application to physical problems) of a previous work communicated in
IMACS-ACA'98. A part of the overall algorithm has been implemented in the
computer algebra system \textsc{MAPLE V}.\\[1cm] \textbf{References:}\\[0.2cm]  1 - G. Cicogna and G. Gaeta, \textit{Poincar\'e normal forms and
Lie point symmetries}, J. Phys. A (Math. Gen.) 27 (1994) 461-476.\\  2 - G.
Cicogna and G. Gaeta, \textit{Normal forms and nonlinear symmetries}, J.
Phys. A (Math. Gen.) 27 (1994) 7115-7124.\\  3 - H. Ito, \textit{Convergence
of Birkhoff normal forms for integrable systems}, Comment. Math. Helv. 64
(1989) 412-461.\\  4 - H. Ito, \textit{Some aspects of integrability and
action-angle variables}, Sugaku Expositions, vol. 3, no. 2 (1990) 213-232.\\ 
5 - H. Ito, \textit{Integrability of Hamiltonian systems and Birkhoff normal
forms in the simple resonance case}, Math. Ann. 292 (1992) 411-444.\\
6 - S. Walcher, \textit{Algebraishe Probleme bei Normalformen Gew\"ohnlicher
Differentialgleichungen}, Preprint Technische Universit\"at M\"unchen, 1990.

\end{document}
